\documentclass[../main.tex]{subfiles}
\begin{document}
\chapter{Introduction}
\begin{remark}[Note]
  Throughout this course, we shall assume that all of our random variables and random processes are defined on an appropriate underlying probability space $(\Omega, \mathcal{F}, \P)$.
\end{remark}
\section{Definitions}
Loosely speaking, a \textit{Markov Chain} is a sequence of random variables where each random variable depends only on the previous random variable in the sequence.
\begin{definition}[Markov Chain]
  A \textit{Markov Chain} is a sequence $\mchain{X} = (X_n)_{n \geq 0}$ of random variables which all take values in the same discrete, countably infinite \textit{state space} $I$ such that:
  \[
    \P(X_{n + 1} = x_{n + 1} \vert X_n = x_n, \ldots, X_0 = x_0) = \P(X_{n + 1} = x_{n + 1} \vert X_n = x_n)
  \]
  for all $n \in \N$ and $x_1, \ldots, x_n \in I$.
\end{definition}
\begin{remark}[Note]
  This does \textbf{not} mean that $X_{n + 1}$ is independent of $X_1, \ldots, X_{n - 1}$ however it \textbf{does} mean that $X_{n + 1}$ is \textit{conditionally independent} from $X_1, \ldots, X_{n - 1}$ \textit{given} that $X_n = x_n$.
  Intuitively, given we know $X_n = x_n$, the values of $X_1, \ldots, X_{n - 1}$ do not contribute to the probabilities for $X_{n + 1}$.
\end{remark}
\begin{definition}[Time-Homogenous Markov Chain]
  We say that a Markov chain is \textit{time-homogeneous} if $\P(X_{n + 1} = y \vert X_n = x)$ is always independent of $n$.
\end{definition}
Intuitively, a Markov chain is time-homogeneous if the probabilities do not depend on what time $n$ we reach a particular state.

In this course, we shall assume that all Markov chains considered are time-homogeneous.
\begin{definition}[Transition Matrix]
  For time-homogeneous Markov chains, we define the transition matrix $P$ as:
  \[
    P_{x y} = P(x, y) = \P(X_{n + 1} = y \vert X_n = x)
  \]
  for all $x, y \in I$.
\end{definition}
\begin{remark}
  We can choose any $n$ for the transition matrix since the Markov chain is time-homogeneous.
\end{remark}
\begin{definition}[Stochastic Matrix]
We call a matrix $P$ \textit{stochastic} if each of its rows are probability distributions on $I$.
That is, $P_{x y} \geq 0\ \forall x, y$ and $\sum_{y \in I} P_{x y} = 1$ for all rows $x \in I$.
\end{definition}
We see that all transition matrices are stochastic because their rows are defined to be conditional distributions given $X_n = x$.
\begin{example}[Binary Markov Chain]
  A simple \textit{binary} Markov Chain has two states $I = \{1, 2\}$ with a probability $\alpha$ to transfer from $1 \to 2$ and probability $\beta$ to go back from $2 \to 1$.

  This has the transition matrix:
  \[
    P = \begin{pmatrix}
    1 - \alpha & \alpha \\
    \beta & 1 - \beta \\
    \end{pmatrix}
  \]
  We can also represent this Markov chain as a diagram:
  \begin{center}
  \begin{tikzpicture}[>=stealth, shorten >=1pt, node distance=2cm, auto, every initial by arrow/.style={dashed}]
    \node[state] (q_0) {$1$};
    \node[state] (q_1) [right=of q_0]{$2$};

    \path[->] (q_0) edge [loop above] node {$1 - \alpha$} ()
                    edge [bend left] node {$\alpha$} (q_1);
    \path[->] (q_1) edge [loop above] node {$1 - \beta$} ()
                    edge [bend left] node {$\beta$} (q_0);
  \end{tikzpicture}
  \end{center}
  The nodes represent the states $I$ of the Markov chain, and the arrows are annotated with the probability of the chain transitioning from the state at the start of the arrow to the state at the end.
\end{example}
\begin{example}
  Consider now a Markov chain with three states $I = \{1, 2, 3\}$ and transition matrix:
  \[
    P = \begin{pmatrix}
    1/3 & 1/3 & 1/3 \\
    1/2 & 0 & 1/2 \\
    0 & 0 & 1 \\
    \end{pmatrix}
  \]
  or in a diagram:
  \begin{center}
  \begin{tikzpicture}[>=stealth, shorten >=1pt, node distance=2cm, auto, every initial by arrow/.style={dashed}]
    \node[state] (q_0) {$1$};
    \node[state] (q_1) [right=of q_0]{$2$};
    \node[state] (q_2) [below=of q_0]{$3$};

    \path[->] (q_0) edge [loop above] node {$\frac{1}{3}$} ()
                    edge [bend left] node {$\frac{1}{3}$} (q_1)
                    edge node {$\frac{1}{3}$} (q_2);
    \path[->] (q_1) edge [bend left] node {$\frac{1}{2}$} (q_0)
                    edge node {$\frac{1}{2}$} (q_2);
    \path[->] (q_2) edge [loop below] node {$1$} ();
  \end{tikzpicture}
  \end{center}
  We see that if we reach state 3, then we get stuck there indefinitely.
\end{example}
\begin{definition}[Initial Distribution]
  We call the vector $\vec{\lambda}$ with components $\lambda_i \in [0, 1]$ for $i \in I$ an \textit{initial distribution} for a Markov Chain $\mchain{X}$ if:
  \[
    \P(X_0 = x) = \lambda_x\ \forall x \in I
  \]
  That is, $\vec{\lambda}$ specifies the probabilities of the Markov chain starting at any of the states $x \in I$.
\end{definition}
\begin{remark}[Notation]
  Due to time-homogeneity, a Markov chain is uniquely specified by a transition matrix $P$ and initial distribution $\vec{\lambda}$ and so we can write $\mchain{X} \sim \markov(\vec{\lambda}, P)$.
\end{remark}
\begin{theorem}
  A sequence of random variables $\mchain{X} = (X_n)_{n \in \N}$ has $\mchain{X} \sim \markov(\vec{\lambda}, P)$ if and only if:
  \[
    \P(X_0 = i_0, X_1 = i_1, \ldots, X_n = i_n) = \lambda_{i_0} P(i_0, i_1) \cdots P(i_{n - 1}, i_{n})
  \]
  for all $n \geq 1$ and paths $i_0, \ldots, i_n \in I$.
\end{theorem}
\begin{proof}
  Continued next lecture.
\end{proof}
\end{document}
